\section{Recovering the format and the rotation basis}
\label{sec:format}

\paragraph{The basis is recoverable.}
Plain absmean round-to-nearest of the rotated public base, with the convention
of Equation~\ref{eq:rotation}, reproduces the released trits at
$0.920$ mean / $0.885$ minimum agreement over $26$ tensors sampled from layers
$0$, $3$, $62$, and $63$ (Table~\ref{tab:gate1}). Chance is $0.33$; a
$1.7$B-era public release of the same family measured $0.612$ under the same
metric, so the $0.92$ is not an artifact of the metric but a property of the
correct basis. Getting there required four conventions to be right
simultaneously, visible as a cliff in the table: the transform is applied on
the input axis, the sign multiplication happens before the Hadamard transform,
the block size is $1024$ with per-block explicit signs, and norms are not
folded into the rotation. Any other combination lands at chance.

\begin{table}[t]
\caption{Per-tensor trit agreement of rotate-and-round hypotheses against the
released codes (mean over tensors; chance $1/3$). Only the input-axis,
sign-then-Hadamard, explicit-signs convention recovers the codes.}
\label{tab:gate1}
\begin{center}
\begin{tabular}{lc}
\toprule
hypothesis & agreement $\eta$ \\
\midrule
no rotation (plain absmean RTN) & $0.333$ \\
output-axis rotation & $\sim 0.33$ \\
FWHT-then-sign order & $\sim 0.33$ \\
sign-then-FWHT, input axis, block $1024$ & $\mathbf{0.920}$ mean / $0.885$ min \\
\bottomrule
\end{tabular}
\end{center}
\end{table}

\paragraph{Layouts matter as much as the basis.}
The released file does not use a single layout for all tensors. The
value-head rows of the linear-attention \code{in\_proj\_qkv} and
\code{in\_proj\_z} matrices are stored tiled rather than grouped; after
applying the standard converter reorder their agreement jumps from $0.56$ and
$0.36$ to $0.927$. The \code{ssm\_out} matrix is the opposite: it stays
grouped, and applying the reorder that helps the other two drops its agreement
from $0.936$ to $0.34$. Feed-forward and attention projections are plain.
These per-suffix layout rules had to be recovered empirically because they
differ from what a current converter would emit.

\paragraph{Scale conventions.}
The stored group scale $d$ is close to least-squares optimal for the released
codes themselves ($d/d_{\mathrm{LS}} \in [1.007, 1.034]$) and about
$1.37\times$ larger than our absmean scale. Substituting their scale into our
rotate-and-round solution \emph{lowers} agreement ($0.88$--$0.91$) relative to
absmean ($0.89$--$0.95$), consistent with the released codes having been chosen
at a decision boundary near $0.5\times$ the absmean scale. The public reference
quantizer exported by the inference fork is a different object: it uses
$d=\max|x|$ and, run on the rotated base, matches only $0.493$ of the released
trits with a scale ratio of about $2.55$, so the fork's public tool is not the
quantizer that produced the release.

\paragraph{Independent confirmation.}
The authors separately publish the format ground truth in a sibling
MLX-oriented repository: a Hadamard manifest (block size, transform name, axis,
explicit sign vectors, and the rotated tensor list) and the runtime codec for
the packed format~\citep{prismbonsai2mlx}. Against that manifest our recovery
matches on every axis (Table~\ref{tab:format-verify}): block size $1024$,
normalized Sylvester Walsh--Hadamard, input-last-dimension, explicit
$\pm 1$ signs, widths $\{5120, 6144, 17408\}$, and a byte-exact match on the
34-byte-per-128-weights codec layout. The manifest lists $401$ forward-rotated
tensors; our census counts $402$ ternary payloads. The difference is the
embedding, which is inverse-rotated and named separately; the two counts
coincide once it is accounted for. This turns the format chapter from
reconstruction into verification: the recovery was derived from the released
file alone, and the manifest confirms it.

\begin{table}[t]
\caption{Our recovered format conventions against the independently published
manifest. All checked conventions match.}
\label{tab:format-verify}
\begin{center}
\begin{tabular}{lll}
\toprule
property & recovered from release & published manifest \\
\midrule
block size & $1024$ & $1024$ \\
transform & normalized Walsh--Hadamard & normalized Sylvester WH \\
axis & input, last dimension & input-last-dimension \\
signs & explicit, per tensor & explicit, $28672$ values \\
widths & $5120 / 6144 / 17408$ & same \\
value-head grouping & \code{ssm\_out} grouped & \code{gdn\_v\_grouped: true} \\
packing & $34$\,B per $128$, $16$ trits/word & byte-exact match \\
forward-rotated set & $401$ tensors & $401$ \code{weight\_names} \\
ternary payloads & $402$ ($+$ inverse embedding) & $401 + $ embedding \\
\bottomrule
\end{tabular}
\end{center}
\end{table}

\paragraph{The zero band is the quantizer, not an attractor.}
Decoding the released file yields the deployed code distribution directly. The
zero fraction is $0.3276$ overall and uniform to three decimals across every
tensor class---attention, feed-forward, linear-attention, output head, and
embedding (Table~\ref{tab:zero}). Sign balance is symmetric ($0.3358$ positive,
$0.3365$ negative), and no invalid code word appears. Under the absmean rule
the expected zero share follows from the shape of the weight distribution
alone: about $0.309$ for a near-Gaussian and about $0.328$ at kurtosis
$\approx 4.5$ (simulated at 4M samples), with our rotated 1.7B canary sitting
at the Gaussian value. The released model sits on the heavy-tailed prediction.
Two consequences follow. First, the sparsity level is not a separately tuned
quantity; there is no free knob for it. Second, the observation is
\emph{consistent with} the zero being the quantizer's central band rather than
a trained attractor, though it does not prove that on its own---a two-well
potential could also land near $1/3$. What it does rule out is any recipe that
needs a tuned sparsity target to reach the released quality.

\begin{table}[t]
\caption{Zero-state structure of the released 27B, decoded from the file.
The zero fraction is uniform across tensor classes.}
\label{tab:zero}
\begin{center}
\begin{tabular}{lrr}
\toprule
class & tensors & zero fraction \\
\midrule
attention $q$ & $64$ & $0.3276$ \\
attention $k$ / $v$ & $32$ & $0.3274$ \\
attention $o$ & $16$ & $0.3275$ \\
feed-forward gate / up / down & $192$ & $0.3276$ \\
linear attention (qkv / z / out) & $96$ & $0.3275$ \\
output head & $1$ & $0.3279$ \\
embedding & $1$ & $0.3280$ \\
\midrule
overall & $402$ & $\mathbf{0.3276}$ \\
\bottomrule
\end{tabular}
\end{center}
\end{table}

\paragraph{Where the missing $8\%$ lives.}
The residual codes were not only fewer; they were structurally different.
Within groups, the released codes are non-monotone in $|w|$ (violations
essentially everywhere), which no pure threshold rule in this basis can
produce. The mismatch against rotate-and-round decomposes into
rotate-and-round weights that the release moved to zero (about $1.6$M elements)
and zeros the release moved off zero (about $1.0$M); there are no
$-1 \leftrightarrow +1$ flips. Mismatches concentrate at $|w| \approx 0.5\times$
the group scale and vanish for $|w| > 1.2\times$ the scale. Measured in weight
space, the released assignment is about $10\%$ closer to the rotated base than
round-to-nearest is (relative error $0.460$--$0.482$ versus
$0.510$--$0.511$). Finally, the tensors exempt from ternarization also moved:
norms drifted by up to $0.07$ in absolute value, and the small BF16
projections in the recurrent path are element-wise decorrelated from the base
while their sorted norms remain correlated at $0.93$--$0.98$, a signature of
permutation plus drift rather than re-initialization.

This fingerprint is consistent with error-compensated quantization or with
training, and it is what motivated the mechanism tests of the next two
sections. The short version is that the first hypothesis did not survive: the
$92\%$ agreement is an excellent container---and a poor proxy for
quality---because the missing codes are exactly where the model's quality
lives.
